\begin{apartat}{1,25}
Un satèl·lit descriu una trajectòria circular de radi $R$ al voltant d'una massa central. El temps que triga a donar-hi una volta sencera és $T$. Deduïu l'expressió per a calcular la intensitat del camp gravitatori, $g$, creat per la massa central en els punts de l'òrbita del satèl·lit en funció dels paràmetres $R$ i $T$. Considereu que la Lluna descriu una òrbita circular al voltant de la Terra amb una distància entre centres de $384\times 10^{6}\ \mathrm{m}$ i amb un període de 27,3 dies. Fent ús només d'aquestes dues dades i de l'expressió trobada anteriorment, calculeu la intensitat del camp gravitatori als punts de l'òrbita de la Lluna.
\end{apartat}

\begin{apartat}{1,25}
Deduïu l'expressió de l'energia cinètica mínima necessària perquè un coet de massa $m$ pugui escapar d'un objecte astronòmic de massa $M$ i radi $R$. Quantes vegades més gran és l'energia cinètica mínima perquè el coet pugui escapar de la Terra respecte de l'energia mínima que necessita per a escapar de la Lluna? (Només podeu fer servir les dades donades tot seguit.)
\end{apartat}

\dades{%
  $M_\mathrm{Terra} = 81,3\ M_\mathrm{Lluna}$.\\
  $R_\mathrm{Terra} = 3,67\ R_\mathrm{Lluna}$.}
